\section*{Mit diesem Übungsheft arbeiten}
Dieses Heft ergänzt die
\href{_B19-cuts-reading.pdf\#cuts.reading}{Lesefassung
\emph{Dedekindsche Schnitte und Vollständigkeit}}.
Acht Aufgaben führen vom Prüfen der Definition zum eigenen
Vollständigkeitsbeweis. Die Lösungen beginnen auf
Seite~\pageref{cuts:solutions}.

\textbf{Die vier Schnittbedingungen.}
Eine Teilmenge $A\subseteq\mathbb Q$ ist ein Schnitt, wenn
\begin{enumerate}
 \item $A$ nicht leer ist,
 \item $A$ nicht ganz $\mathbb Q$ ist,
 \item mit $p\in A$ auch jedes rationale $q<p$ in $A$ liegt,
 \item zu jedem $p\in A$ ein $r\in A$ mit $p<r$ existiert.
\end{enumerate}
Die letzten beiden Bedingungen heißen \emph{abwärts abgeschlossen}
und \emph{ohne größtes Element}.

Für $q\in\mathbb Q$ schreiben wir
\[
 \widehat q=\{p\in\mathbb Q:p<q\}.
\]
Die Menge aller Schnitte heißt $\mathbb R$. Ihre Ordnung ist die
Inklusion: $A\le B$ bedeutet $A\subseteq B$.
Ein Element von $A$ ist also eine rationale Zahl;
ein Element einer Schnittfamilie $\mathcal A$ ist ein ganzer Schnitt.

\textbf{So arbeitest du.}
Zeichne zuerst eine Skizze oder suche einen einzelnen Zeugen.
Für $A\subseteq B$ nimm ein beliebiges Element von $A$;
für $A\nsubseteq B$ genügt ein Element von $A\setminus B$.
Schreibe bei jeder Vereinigung ausdrücklich auf, aus welchem
Familienglied ein betrachtetes Element stammt.

Die Karofläche rechts gehört jeweils zur Textseite links.
Nutze sie auch für einen zweiten Beweisversuch. Die Aufgaben bauen
aufeinander auf; frühere Ergebnisse dürfen später verwendet werden.
Die Aufgaben~7 und~8 erklären, wozu die Voraussetzungen des
Vollständigkeitssatzes nötig sind.


